\section{Conclusion}\label{sec:conclusion}
Neural Bellman Operators organize economic computation around the continuation object that is evaluated, the feasible action that it changes, and the economic consequence of that change. The separation matters for recursive utility, preference adjustment, temporal selves, and strategic interaction because each problem has its own evaluation equation and deviation criterion. The approximation and verification arguments make those distinctions explicit.

The capital economy permits an exact finite Bellman account under the candidate's occupation law, a computable feasible-actor gap, and a quantified transfer to the original continuous diffusion. These results make the evaluation--action--payoff chain executable without identifying a finite-rollout derivative with a continuous-time costate. The independent paired transfer also makes use of the actual common innovations in direct method comparisons, while preserving the original observations and confidence event.

The complete sixteen-stream experiment establishes economically material NBO payoff gains over the declared neural HJB implementation in both ten and fifty dimensions. In dimension fifty, NBO reaches the prescribed online gain target in every stream with less measured complete work than direct policy optimization. Cached Raw remains cheaper; the direct NBO--Raw and NBO--DPO payoff intervals do not resolve their ranking or establish practical equivalence. Thus the experiment identifies a useful domain for the NBO evaluation--actor construction under an explicit computational contract.

The separate mechanism calculation reports all occupation, action, sampling, and numerical allowances, but its sufficient lower bound does not certify positive mechanism welfare at the executed precision. The full-class regret bounds likewise remain wider than the direct-comparison margin. The independent scalar calculation provides a sharper finite-grid account with a classical training and deployment comparison, and preserves the observed state-dependent ranking reversal. These qualifications attach each accuracy statement to the economic and numerical object that it actually controls.

Finite observations also support an operational use of the computed policy. A manager receives measured capital and selects consumption at stated dates; after the allowance for that grid and measurement precision, six of twelve specifications retain a positive continuous-economy lower gain. The original recursive-utility, endogenous-preference, temporal-self, and game applications retain their complete models and proofs. Together they state the economic conditions under which an evaluated continuation object can guide feasible decisions, while the capital experiments quantify the resulting payoff and computational tradeoffs.
